% Source: https://artofproblemsolving.com/community/c3h1984801_help_me_please
% Stable identity: aops-topic-1984801-q01
The parabola meets the $x$-axis at $x=-1$ and $x=5$, so $A=(-1,0)$ and $D=(5,0)$ (with $A$ to the left). The height is $2$ and $\angle BAD=60^\circ$, so $B$ is $2$ above the axis and $\dfrac{2}{\tan60^\circ}=\dfrac{2}{\sqrt3}$ to the right of $A$:
\[
B=\left(-1+\dfrac{2}{\sqrt3},\ 2\right).
\]
By symmetry about the axis $x=2$ of the parabola, $C=\left(5-\dfrac{2}{\sqrt3},2\right)$ is then also on the parabola.

\begin{asy}
usepackage("amsmath");
size(6cm);
defaultpen(fontsize(10pt));
import graph;
real a=-(3+9*sqrt(3))/52;
real f(real x){return a*(x+1)*(x-5);}
draw(graph(f,-1.6,5.6),gray(0.4));
draw((-2,0)--(6.2,0),Arrow(TeXHead));draw((0,-1)--(0,4),Arrow(TeXHead));
pair A=(-1,0),D=(5,0),B=(-1+2/sqrt(3),2),C=(5-2/sqrt(3),2);
draw(A--B--C--D--cycle,linewidth(1));
draw(B--(B.x,0),dashed);label("$2$",(B.x,1),E);
draw(arc(A,0.45,0,60));label("$60^\circ$",A+(0.6,0.25),fontsize(6pt));
dot(A);dot(B);dot(C);dot(D);
label("$A$",A,SW);label("$D$",D,SE);label("$B$",B,NW);label("$C$",C,NE);label("$x$",(6.2,0),S);label("$y$",(0,4),W);
label("\mbox{}",(min(currentpicture,user=true)+max(currentpicture,user=true))/2,black); // keeps the text after the figure black when LaTeX typesets the labels inline
\end{asy}

Substituting $B$:
\[
2=a\cdot\dfrac{2}{\sqrt3}\cdot\left(\dfrac{2}{\sqrt3}-6\right)=a\left(\dfrac{4}{3}-4\sqrt3\right),
\]
so
\[
a=\dfrac{2}{\dfrac{4}{3}-4\sqrt3}=\dfrac{3}{2-6\sqrt3}=\dfrac{3(2+6\sqrt3)}{4-108}=\boxed{-\dfrac{3+9\sqrt3}{52}}\approx-0.357 .
\]
